\section{Knowledge Graphs}

Knowledge graphs (KGs) are structured representations of knowledge that organise information as a graph consisting of entities and the semantic relationships between them. Formally, a knowledge graph can be defined as a directed labelled graph \(G=(V,E)\), where the vertices \(V\) represent entities or concepts, and the edges \(E\) represent semantic relationships between those entities. Unlike relational databases, which primarily model tabular data, knowledge graphs explicitly capture the interconnected nature of information, thereby enabling reasoning over relationships rather than isolated data items.

A common representation of knowledge graphs is the Resource Description Framework (RDF), in which knowledge is expressed as triples of the form \textit{subject--predicate--object}. For example, the statement \emph{``Derivative'' is a prerequisite of ``Integration''} can be represented as the triple \texttt{(Derivative, prerequisiteOf, Integration)}. By combining a large collection of such triples, knowledge graphs provide a structured semantic representation that can be traversed, queried, and extended.

Knowledge graphs have been widely adopted in numerous application domains, including semantic search, question answering, recommendation systems, biomedical informatics, and digital libraries. Their ability to explicitly model semantic relationships enables systems to discover indirect connections between concepts, support multi-hop reasoning, and improve information retrieval beyond simple keyword or embedding similarity. Consequently, knowledge graphs have become an important component of many modern artificial intelligence systems.

In RAG, knowledge graphs provide an alternative representation of the underlying corpus. Instead of retrieving documents solely according to vector similarity, graph-based retrieval methods exploit the semantic relationships between entities to identify relevant concepts and their neighbouring nodes. This structured representation allows retrieval algorithms to incorporate contextual information that may not be captured by dense embeddings alone, thereby improving the completeness and coherence of retrieved knowledge.


The proposed framework adopts this principle by representing the MathHub repository as a structured knowledge graph. Rather than relying exclusively on semantic similarity between documents, the retrieval process exploits prerequisite relationships among mathematical concepts together with learner-specific information to identify educational resources that are both semantically relevant and appropriate for the learner's current level of understanding. This graph-based representation provides the foundation for the GraphRAG-inspired retrieval strategy presented in the following sections.
